\subsection{Beginning Comment}\label{sec:BeginningComment}
A Java project will not consist from Java source files only; aside the already mentioned module definition files(see \tqref{sec:ModuleDefinition}\index{module-info.java}), the package documentation files\index{package-info.java} (refer to \tqfullvref{sec:PackageDocumentation}), there can be various build files, scripts, resource and property files, and even source files other programming languages.

All source files of a project – not only the Java sources – have to start with a comment that gives the copyright\index{copyright} notice and the license\index{license} information\footnote{This is different from the Sun coding conventions\autocite{SUN_CODE_CONVENTIONS:BeginningComments} that recommends to put also the class name and the date into the beginning comment. I omitted the name of the class because this is obvious from the file name in case of a \lstinline|public| class, and it would confuse the reader if the file contains more than one class. The date is considered obsolete.}, if possible at all: it should be obvious that it is sometimes difficult to add this kind of information to a binary resource like an image or a sound file.

For a Java file (a file with a name ending on \verb#*.java#), that beginning comment looks like this:
\begin{lstlisting}[numbers=left,caption={Beginning Comment for a Java file}]
/*
 * ==================================================================
 * Copyright © <Year> <Copyright Notice>
 * ==================================================================
 *
 * <License Notice>
 */
\end{lstlisting}

The same comment block can be used for Groovy\index{Groovy} source files, Kotlin\index{Kotlin} source files, Gradle\index{Gradle} build files (\lstinline|build.gradle|\index{build.gradle}) and generally for all file formats that uses the pattern \lstinline|/* … */| for comments.

A \textit{resource} file\index{resource file}\autocite{ORACLE_DOC_RESOURCEBUNDLE_CLASS} can have XML format or it may contain simple key-value pairs. In the latter form, comments are prepended with the hash symbol (``\verb|#|''), and the beginning comment has the following form: 
\begin{lstlisting}[numbers=left,language=,caption={Beginning Comment for a Resource or a Script file}]
#
# ===================================================================
# Copyright © <Year> <Copyright Notice>
# ===================================================================
#
# <License Notice>
#
\end{lstlisting}

Obviously, this can be used for all other file formats that use the same comment style, like Unix shell scripts\index{shell script}.

For a resource file in XML\index{XML} format\footnote{For an XML file, the beginning comments follows the initial \lstinline|<?xml …?>|, so it is not starting on line~1. That's similar for HTMl files.} (or any other XML file that is part of the sources), as well as for HTML\index{HTML} files, the beginning comment would look like this:
\begin{lstlisting}[numbers=left,firstnumber=5,language=XML,caption={Beginning Comment for XML and HTML files}]
<!--
=====================================================================
 Copyright © <Year> <Copyright Notice>
=====================================================================

 <License Notice>
-->
\end{lstlisting}

\TeX/\LaTeX\index{TeX}\index{LaTeX} uses the per cent sign (``\verb#%#'') to indicate a comment. So the beginning comment for a \TeX/\LaTeX\index{TeX}\index{LaTeX} source file would be like this:
\begin{lstlisting}[numbers=left,language=,caption={Beginning Comment Resource file}]
%
% ===================================================================
% Copyright © <Year> <Copyright Notice>
% ===================================================================
%
% <License Notice>
%
\end{lstlisting}

\verb#<Year>#, \verb#<Copyright Notice># and \verb#<License Notice># has to be replaced by the appropriate texts. The special character \bdq{\copyright} is used not only because is looks better; it tells you immediately whether your environment is capable to deal with UTF-8 encoded files.

The lines with the equals signs (“===”) will end at column~80; refer to chapter \tqfullvref{sec:LineLength}.\footnote{Keep in mind that in this document the line length is set to 70~columns; this means that you cannot just copy and paste the comment blocks into your development environment.}

\subsubsection{Principles}
\begin{enumerate}[label=P\arabic*.]
\item{Each file that is part of the sources for the program/application to build should have a beginning comment.}
\item{That beginning comment must contain the copyright notice with the year and the copytight holder.}
\item{It should contain a license notice (either the license itself or a reference to it).}
\item{No other information should be added to this comment block.}
\end{enumerate}

%==============================================================================
% End of file!!
%==============================================================================
